% ==============================================================================
% Estilo de los teoremas y secciones
% ==============================================================================
\usepackage{titlesec}

% Ajusta el formato de \section: {formato}{etiqueta}{separación}{código antes}{código después}
% Estructura: \titlespacing*{comando}{margen izquierdo}{espacio antes}{espacio después}

% --- SECCIÓN: Centrada, grande y numerada ---
\titleformat{\section}[block]
{\normalfont\Large\bfseries\color{black}\centering}
{\thesection.}
{0.5em}
{}
\titlespacing*{\section}{0pt}{*2}{*3}

% --- SUBSECCIÓN: A la izquierda, normal y SIN numeración ---
\titleformat{s}
{\normalfont\normalsize\bfseries\color{black}\raggedright}
{}
{0em}
{}
\titlespacing*{\subsection}{0pt}{*2}{*1}

% --- SUBSUBSECCIÓN: Más grande (ej. \large) y a la izquierda ---
\titleformat{\subsubsection}
{\normalfont\large\bfseries\color{black}\raggedright} % Cambia \large por \Large si la quieres aún más grande
{\thesubsubsection} % Pon {} aquí si tampoco quieres que esté numerada
{0.5em}
{}
\titlespacing*{\subsubsection}{0pt}{*1.5}{*1}

% --- 2. CONFIGURACIÓN DE LAS SUBSECCIONES 

\newenvironment{enunciado}
  {\par\large\color{black}}
  {\par\medskip}

% Entorno para la resolución (Azul)
\newenvironment{solucion}
  {\par\large\color{blue}}
  {\par\bigskip}

\newtheoremstyle{Proposición}
  {5pt}         % Espacio arriba
  {10pt}        % Espacio abajo
  {\normalfont} % Fuente del cuerpo
  {}            % Sangría del cuerpo 
  {\bfseries}   % Fuente del encabezado (Definición 3.1)
  {.}           % Puntuación después del encabezado
  {.5em}           % Espacio después del encabezado
  {\thmname{#1}\thmnumber{ #2}\thmnote{ (#3)}} % Formato del encabezado

\newtheoremstyle{Ejemplo}
  {5pt}         % Espacio arriba
  {10pt}        % Espacio abajo
  {\normalfont} % Fuente del cuerpo
  {}            % Sangría del cuerpo 
  {\bfseries}   % Fuente del encabezado (Definición 3.1)
  {.}           % Puntuación después del encabezado
  {.5em}           % Espacio después del encabezado
  {Ejemplo \thmnumber{#2}} % Formato del encabezado

\newtheoremstyle{Ejercicio}
  {5pt}         % Espacio arriba
  {10pt}        % Espacio abajo
  {\normalfont} % Fuente del cuerpo
  {0pt}            % Sangría del cuerpo 
  {\bfseries}   % Fuente del encabezado 
  {}           % Puntuación después del encabezado
  {0pt}           % Espacio después del encabezado
  {Ejercicio \thmnumber{#2}\newline} % Formato del encabezado

\newtheoremstyle{Ampliacion}
  {5pt}         % Espacio arriba
  {10pt}        % Espacio abajo
  {\normalfont} % Fuente del cuerpo
  {0pt}            % Sangría del cuerpo 
  {\bfseries}   % Fuente del encabezado 
  {}           % Puntuación después del encabezado
  {0pt}           % Espacio después del encabezado
  {\centering Complemento \newline} % Formato del encabezado


%-----------------------------------------------------------------------------------------------
%-----------------------------------------------------------------------------------------------

\theoremstyle{Proposición}
\newtheorem*{prop}{Proposición}

\tcolorboxenvironment{prop}{
  enhanced,
  breakable,                % Permite que el ejemplo salte de página si es muy largo
  colback=black!5!white,    % Fondo verde muy claro (5% verde, 95% blanco)
  colframe=black!50!black,  % Borde verde oscuro
  boxrule=0.5pt,            % Grosor del borde
  arc=3pt,                  % Esquinas redondeadas
  top=2pt, bottom=2pt,      % Espaciado interior arriba y abajo
  left=5pt, right=5pt       % Espaciado interior a los lados
}

\newtheorem*{deff}{Definición}
\tcolorboxenvironment{deff}{
  enhanced,
  breakable,                % Permite que el ejemplo salte de página si es muy largo
  colback=black!5!white,    % Fondo verde muy claro (5% verde, 95% blanco)
  colframe=black!50!black,  % Borde verde oscuro
  boxrule=0.5pt,            % Grosor del borde
  arc=3pt,                  % Esquinas redondeadas
  top=2pt, bottom=2pt,      % Espaciado interior arriba y abajo
  left=5pt, right=5pt       % Espaciado interior a los lados
}

\theoremstyle{Ejemplo}
\newtheorem{ej}{Ejemplo}
\tcolorboxenvironment{ej}{
  enhanced,
  breakable,                % Permite que el ejemplo salte de página si es muy largo
  colback=green!5!white,    % Fondo verde muy claro (5% verde, 95% blanco)
  colframe=green!50!black,  % Borde verde oscuro
  boxrule=0.5pt,            % Grosor del borde
  arc=3pt,                  % Esquinas redondeadas
  top=2pt, bottom=2pt,      % Espaciado interior arriba y abajo
  left=5pt, right=5pt       % Espaciado interior a los lados
}

\theoremstyle{Ejercicio}
\newtheorem{ex}{}
\tcolorboxenvironment{ex}{
  enhanced,
  breakable,                % Permite que el ejemplo salte de página si es muy largo
  colback=red!5!white,    % Fondo verde muy claro (5% verde, 95% blanco)
  colframe=red!50!black,  % Borde verde oscuro
  boxrule=0.5pt,            % Grosor del borde
  arc=3pt,                  % Esquinas redondeadas
  top=2pt, bottom=2pt,      % Espaciado interior arriba y abajo
  left=5pt, right=5pt       % Espaciado interior a los lados
}


\theoremstyle{Ejemplo}
\newtheorem{_ejs}[ej]{Ejemplo}
\newenvironment{ejs}
  {\begin{_ejs}\leavevmode\vspace{-0.2\baselineskip}\begin{itemize}[leftmargin=15pt]}
  {\end{itemize}\end{_ejs}}

\tcolorboxenvironment{ejs}{
  enhanced,
  breakable,                % Permite que el ejemplo salte de página si es muy largo
  colback=green!5!white,    % Fondo verde muy claro (5% verde, 95% blanco)
  colframe=green!50!black,  % Borde verde oscuro
  boxrule=0.5pt,            % Grosor del borde
  arc=3pt,                  % Esquinas redondeadas
  top=2pt, bottom=2pt,      % Espaciado interior arriba y abajo
  left=5pt, right=5pt       % Espaciado interior a los lados
}


\newenvironment{amp}
  {\begin{center}
  \large\textbf{Ampliación}
  \end{center}}
  {}

\tcolorboxenvironment{amp}{
  enhanced,
  breakable,                % Permite que el ejemplo salte de página si es muy largo
  colback=blue!5!white,    % Fondo verde muy claro (5% verde, 95% blanco)
  colframe=blue!50!black,  % Borde verde oscuro
  boxrule=0.5pt,            % Grosor del borde
  arc=3pt,                  % Esquinas redondeadas
  top=2pt, bottom=2pt,      % Espaciado interior arriba y abajo
  left=5pt, right=5pt       % Espaciado interior a los lados
}